\documentclass[12pt,a4paper]{article}
\usepackage[utf8]{inputenc}
\usepackage[T1]{fontenc}
\usepackage[spanish]{babel}
\usepackage{geometry}
\usepackage{hyperref}
\usepackage{booktabs}
\usepackage{xcolor}
\usepackage{longtable}
\usepackage{fancyhdr}
\usepackage{titlesec}
\usepackage{array}
\usepackage{parskip}

\geometry{
  top=2.5cm,
  bottom=2.5cm,
  left=3cm,
  right=3cm
}

\definecolor{criticalred}{RGB}{180,30,30}
\definecolor{warnorange}{RGB}{200,100,0}
\definecolor{okgreen}{RGB}{30,130,60}
\definecolor{accentblue}{RGB}{30,80,160}
\definecolor{lightgray}{RGB}{245,245,245}
\definecolor{titlegray}{RGB}{60,60,80}

\hypersetup{
  colorlinks=true,
  linkcolor=accentblue,
  urlcolor=accentblue,
  pdftitle={forge v2.0 — Analisis Dual + Benchmark},
  pdfauthor={Analisis independiente},
  pdfsubject={Framework de gobernanza de agentes IA}
}

\pagestyle{fancy}
\fancyhf{}
\fancyhead[L]{\small\textcolor{titlegray}{forge v2.0 — Analisis Dual + Benchmark (v4)}}
\fancyhead[R]{\small\textcolor{titlegray}{Mayo 2026}}
\fancyfoot[C]{\small\thepage}
\renewcommand{\headrulewidth}{0.4pt}

\titleformat{\section}
  {\large\bfseries\color{titlegray}}
  {\thesection}{1em}{}[\titlerule]

\titleformat{\subsection}
  {\normalsize\bfseries\color{accentblue}}
  {\thesubsection}{1em}{}

\titleformat{\subsubsection}
  {\normalsize\bfseries}
  {\thesubsubsection}{1em}{}

\newcolumntype{L}[1]{>{\raggedright\arraybackslash}p{#1}}
\newcolumntype{C}[1]{>{\centering\arraybackslash}p{#1}}

\begin{document}

%% ============================================================
%% PORTADA
%% ============================================================
\begin{titlepage}
  \centering
  \vspace*{2cm}
  {\fontsize{28}{34}\selectfont\bfseries\color{titlegray} forge v2.0}\par
  \vspace{0.5cm}
  {\fontsize{18}{22}\selectfont\color{accentblue} Analisis Dual + Benchmark Competitivo}\par
  \vspace{0.3cm}
  {\large\color{titlegray} Version 4 del ciclo de analisis}\par
  \vspace{2cm}
  \hrule height 1pt
  \vspace{0.5cm}

  \begin{abstract}
    \noindent
    Este documento sintetiza dos analisis independientes realizados sobre
    \textbf{forge v2.0} (44 commits, un unico maintainer, rama \texttt{main}).
    El primer analisis audito el codigo fuente en busca de bugs, inconsistencias
    entre documentacion y codigo, y riesgos de adopcion. El segundo comparo forge
    con cinco alternativas del ecosistema de desarrollo asistido por IA
    (aider, Cursor rules, cline/Roo Code, OpenHands, DIY manual) en diez criterios
    cuantificados.

    Los hallazgos incluyen cuatro bugs confirmados en el codigo ---el mas critico
    implica que la integracion CI documentada genera falsa seguridad--- y una
    ventaja competitiva real: forge lidera en 8 de 10 criterios del benchmark
    en su nicho especifico de gobernanza de equipos de agentes IA.

    La sintesis presenta un veredicto diferenciado por perfil de equipo.
  \end{abstract}

  \vspace{0.5cm}
  \hrule height 1pt
  \vspace{2cm}

  \begin{tabular}{rl}
    \textbf{Fecha:} & 2026-05-03 \\
    \textbf{Version analizada:} & forge v2.0, 44 commits \\
    \textbf{Metodologia:} & Analisis dual independiente \\
    \textbf{Benchmark:} & 5 alternativas, 10 criterios \\
  \end{tabular}

  \vfill
  {\small\color{titlegray} Analisis independiente. No afiliado con Anthropic ni con
  ningun vendor mencionado en el benchmark.}
\end{titlepage}

%% ============================================================
%% TABLA DE CONTENIDOS
%% ============================================================
\tableofcontents
\newpage

%% ============================================================
\section{Resumen ejecutivo}
%% ============================================================

forge es un framework de gobernanza de agentes IA que define, versiona y despliega
equipos de agentes especializados en multiples runtimes (Claude Code, OpenCode, Kiro)
desde una unica fuente de verdad declarativa (\texttt{project.yaml}).

En su version 2.0, post-mejoras, el proyecto ha alcanzado un nivel de madurez notable:
358 tests automatizados, 13 profiles de stack, un CLI interactivo sin dependencias
externas y soporte declarativo para compliance regulatorio (GDPR, Ley 21.719 chilena).

Al mismo tiempo, cuatro problemas estructurales persisten. Uno de ellos ---la ausencia
del campo \texttt{summary} en el JSON de auditoria--- es un bug que hace que la integracion
CI documentada en el README oficial genere falsa seguridad en produccion.

El analisis benchmark concluye que, para equipos de 2 a 8 personas que usan Claude Code
de forma activa con multiples proyectos, forge no tiene competidor directo en el ecosistema.
Para otros perfiles, las alternativas son mas adecuadas.

%% ============================================================
\section{Estado del proyecto}
%% ============================================================

\begin{center}
\begin{tabular}{ll}
  \toprule
  \textbf{Metrica} & \textbf{Valor} \\
  \midrule
  Commits totales & 44 \\
  Autores & 1 (Cristian Correa) \\
  PRs externos & 0 visibles \\
  Tests automatizados & 358 (2.86 segundos) \\
  Profiles de stack (Tier 2) & 13 \\
  Agentes core (Tier 1) & 7 \\
  MCP servers en catalogo & 20 \\
  Runtimes soportados & 3 (Claude Code, OpenCode, Kiro) \\
  \bottomrule
\end{tabular}
\end{center}

La arquitectura central es un sistema de tres tiers: Tier 1 (agentes universales en
\texttt{core/agents/}), Tier 2 (profiles por stack en \texttt{profiles/<stack>/}) y
Tier 3 (agentes de dominio en el proyecto del equipo, fuera de forge). La regla de
colision garantiza que la especializacion por stack no rompe el roster base.

%% ============================================================
\section{Bugs confirmados en el codigo}
%% ============================================================

Los siguientes son hechos verificados entre el codigo fuente y la documentacion,
con referencias precisas a archivos y numeros de linea. No son estimaciones de riesgo.

\subsection{Bug 1 — Campo \texttt{summary} ausente en el JSON de auditoria}

\textbf{\textcolor{criticalred}{Severidad: ALTA}}

El comando \texttt{forge-audit.py --json} emite un JSON con cuatro campos:
\texttt{project}, \texttt{agents}, \texttt{opportunities} y \texttt{orphans}.
El campo \texttt{summary} no existe en el output.

Sin embargo, los siguientes archivos documentan la integracion CI usando ese campo
inexistente:

\begin{itemize}
  \item \texttt{README.md}, linea 173: \texttt{jq -e '.summary.errors == 0'}
  \item \texttt{docs/guide.md}, linea 332: \texttt{jq -e '.summary.errors == 0'}
  \item \texttt{forge.py}, lineas 696 y 967: \texttt{jq '.summary'} y
        \texttt{jq '.summary.errors'}
\end{itemize}

En \texttt{jq}, acceder a un campo inexistente devuelve \texttt{null} sin error y
sin exit code distinto de 0. Un pipeline de CI construido con los ejemplos del README
nunca fallara aunque haya errores criticos de auditoria. Ningun test en
\texttt{test\_forge\_audit.py} verifica la estructura del JSON de salida.

\textbf{Impacto:} Falsa seguridad en CI. Un equipo que siga la documentacion oficial
tendra un pipeline que reporta exito ante cualquier estado de los agentes.

\subsection{Bug 2 — Flag \texttt{--forge} documentada pero no implementada}

\textbf{\textcolor{warnorange}{Severidad: MEDIA}}

\texttt{docs/guide.md} referencia el flag \texttt{--forge .agentic} en cinco ocasiones
(lineas 112, 146, 180, 205, 282). El script \texttt{forge-audit.py} no implementa ese
parametro. Al ejecutar el comando documentado, el flag se ignora silenciosamente: no
hay error, no hay advertencia, el comportamiento es diferente al documentado.

\subsection{Bug 3 — Opcion \texttt{--only} del menu no implementada}

\textbf{\textcolor{warnorange}{Severidad: MEDIA}}

\texttt{forge.py}, linea 712, invoca \texttt{forge-audit.py --only=\{agent\}}.
El script no implementa \texttt{--only}: ejecuta siempre el audit completo. La opcion
del menu describe auditar ``sin revisar el roster completo'' y hace lo contrario.

\subsection{Bug 4 — Documentacion desactualizada en tres archivos}

\textbf{Severidad: BAJA}

\texttt{forge.py} (linea 754) describe ``9 profiles'' cuando el repositorio tiene 13.
\texttt{README.md} no lista astro, django, vuenuxt, go-gin ni sveltekit en la tabla
Tier 2. \texttt{templates/project.yaml.tpl} no incluye los cuatro profiles nuevos.
La causa: la documentacion no se actualizo al agregar los profiles en v2.0.

%% ============================================================
\section{Benchmark comparativo}
%% ============================================================

\subsection{Herramientas evaluadas}

\begin{description}
  \item[aider] (44.3k estrellas) CLI de pair programming. Soporta 100+ lenguajes,
    commits automaticos, repomap para contexto. Sin profiles, sin gobernanza de equipo,
    sin auditoria.
  \item[Cursor rules] Sistema de reglas del editor Cursor. Archivos de contexto por
    directorio. Simple y efectivo para un desarrollador individual. No multi-runtime,
    no profiles, no auditoria.
  \item[cline / Roo Code] (61.3k / 23.8k estrellas) Extensiones de VS Code con agente
    autonomo. Multiples LLM providers, MCP, checkpoints. Sin profiles de stack, sin
    fuente de verdad declarativa.
  \item[OpenHands] (72.6k estrellas) Plataforma de agentes autonomos. SDK Python, CLI,
    GUI. Sin profiles predefinidos, orientado a tareas autonomas, no a gobernanza.
  \item[DIY manual] Crear \texttt{.claude/agents/} a mano sin framework. Maxima
    flexibilidad, entropia creciente con el tiempo.
\end{description}

\subsection{Matriz de benchmark}

\begin{longtable}{L{4.5cm}C{1.2cm}C{1.2cm}C{1.4cm}C{1.4cm}C{1.4cm}C{1.0cm}}
  \caption{Benchmark forge vs. alternativas (escala 1--5)} \\
  \toprule
  \textbf{Criterio} & \textbf{forge} & \textbf{aider} & \textbf{Cursor} &
  \textbf{cline/Roo} & \textbf{OpenHands} & \textbf{DIY} \\
  \midrule
  \endfirsthead
  \toprule
  \textbf{Criterio} & \textbf{forge} & \textbf{aider} & \textbf{Cursor} &
  \textbf{cline/Roo} & \textbf{OpenHands} & \textbf{DIY} \\
  \midrule
  \endhead
  \midrule
  \multicolumn{7}{r}{\small Continua en la pagina siguiente\ldots} \\
  \endfoot
  \bottomrule
  \endlastfoot
  Setup $<$ 10 min        & \textbf{5} & 5 & 4 & 4 & 3 & 1 \\
  Especializacion por stack & \textbf{5} & 1 & 2 & 2 & 1 & 3 \\
  Multi-runtime             & \textbf{5} & 1 & 1 & 1 & 2 & 2 \\
  CLI interactivo           & \textbf{5} & 4 & 1 & 3 & 4 & 1 \\
  Auditoria de agentes      & \textbf{5} & 1 & 1 & 1 & 1 & 1 \\
  Catalogo MCP              & \textbf{5} & 1 & 1 & 3 & 2 & 1 \\
  Tests del framework       & \textbf{5} & 3 & 1 & 3 & 4 & 1 \\
  Sin vendor lock-in        & 5 & 5 & 1 & 4 & 4 & \textbf{5} \\
  Comunidad activa          & 2 & 5 & 4 & \textbf{5} & 5 & -- \\
  Gobernanza y compliance   & \textbf{5} & 1 & 1 & 1 & 2 & 2 \\
  \midrule
  \textbf{Total}            & \textbf{47} & 27 & 17 & 27 & 29 & 21 \\
\end{longtable}

forge lidera en 8 de 10 criterios. Las dos excepciones son significativas:
\textbf{comunidad activa} (donde las alternativas tienen comunidades 10--50 veces
mayores) y \textbf{vendor lock-in real} (donde el informe critico identifico que
\texttt{orchestrator.md} usa APIs exclusivas de Claude Code).

\subsection{Interpretacion del benchmark}

Las herramientas comparadas no resuelven el mismo problema. aider es una herramienta
de edicion asistida sin estado entre proyectos. Cursor rules es configuracion por
directorio. cline y Roo Code dependen de VS Code. OpenHands es una plataforma de
ejecucion autonoma a escala.

La ventaja diferencial de forge es concreta: 13 profiles con instrucciones especializadas
por stack que nadie en el ecosistema provee empaquetadas. Un \texttt{admin-engineer}
que prohíbe explicitamente Tailwind 3 y SWR, que exige WCAG 2.1 AA, que tiene scope
de archivos definido ---eso no se consigue con ninguna alternativa sin escribirlo a
mano en cada proyecto.

%% ============================================================
\section{Fortalezas}
%% ============================================================

\textbf{Arquitectura de tres tiers.} El sistema Tier 1 / Tier 2 / Tier 3 con regla
de colision precisa garantiza que la especializacion por stack no rompe el roster base.
\texttt{forge-init.py} instala profiles primero, core despues, sin sobreescribir.

\textbf{Auditoria como ciudadano de primera clase.} \texttt{forge-audit.py} detecta
frontmatter incompleto, modelo incorrecto por tier, similitud fuera de umbrales
(\texttt{SIMILARITY\_WARN=0.80}, \texttt{SIMILARITY\_OUTDATED=0.50}), agentes huerfanos
y profiles no usados. El bug del campo \texttt{summary} no invalida la capacidad de
auditoria; invalida la integracion CI documentada.

\textbf{358 tests en 2.86 segundos.} Cobertura de auditoria, wizard, integracion
completa de init, adapters, teardown, profiles y CLI. Inusual para un framework de
este tipo.

\textbf{CLI sin dependencias.} TUI completa (pills de categoria, bordes redondeados,
cursor con seleccion, panel de descripcion contextual) en Python puro.

\textbf{Compliance propagado.} El campo \texttt{compliance.frameworks} en
\texttt{project.yaml} activa el \texttt{compliance-reviewer} con modelo \texttt{opus}
obligatorio y propaga las reglas al steering de Kiro. Sin equivalente en el ecosistema.

%% ============================================================
\section{Limitaciones}
%% ============================================================

\textbf{Lock-in con Claude Code a nivel de orchestrator.} El agente central usa
\texttt{Agent()}, \texttt{subagent\_type}, \texttt{run\_in\_background},
\texttt{SendMessage} e \texttt{isolation: worktree} ---APIs exclusivas de Claude Code.
Los adapters de OpenCode y Kiro generan configuraciones de roster pero no resuelven
la brecha de comportamiento. Un equipo en OpenCode no accede al mecanismo de
coordinacion que el orchestrator asume.

\textbf{Bug de CI en produccion.} El campo \texttt{summary} no existe en el JSON de
\texttt{forge-audit.py}. Cualquier pipeline construido con los ejemplos del README
genera falsa seguridad. Es el problema mas urgente.

\textbf{Modelo de instalacion por submodule.} Sin versioning semantico, sin releases
etiquetados, sin script de actualizacion automatizado. El proceso de actualizacion
requiere 6 pasos manuales. \texttt{forge-teardown.py} ignora silenciosamente errores
si el submodule ya fue removido manualmente.

\textbf{Bus factor 1.} 44 commits, un unico autor, cero PRs externos visibles. Para
un framework que gestiona configuracion de agentes en produccion, esto es un riesgo
operativo real para adopcion enterprise.

\textbf{Tests estructurales, no de flujo real.} Los 358 tests verifican que archivos
se copian, que el frontmatter existe, que el YAML contiene ciertos strings. No hay
tests end-to-end del flujo wizard $\rightarrow$ init $\rightarrow$ audit $\rightarrow$
update.

\textbf{Stacks sin profile.} El wizard ofrece Laravel y Angular. No existen profiles
para ninguno. El wizard termina con \texttt{profiles: []} despues de 10 pantallas.

%% ============================================================
\section{Matriz de decision}
%% ============================================================

\begin{center}
\begin{tabular}{L{5cm}L{3cm}L{5.5cm}}
  \toprule
  \textbf{Perfil} & \textbf{Recomendacion} & \textbf{Razon} \\
  \midrule
  Equipo 2--8 personas, Claude Code activo, multiples proyectos &
    \textcolor{okgreen}{\textbf{Adoptar}} &
    Maximo beneficio de profiles y auditoria \\
  \addlinespace
  Equipo con compliance GDPR / Ley 21.719 &
    \textcolor{okgreen}{\textbf{Adoptar}} &
    Sin equivalente en el ecosistema \\
  \addlinespace
  Desarrollador individual, un proyecto, sin compliance &
    \textcolor{warnorange}{DIY o Cursor} &
    Overhead no justificado \\
  \addlinespace
  Migracion planificada a otro runtime &
    \textcolor{warnorange}{Esperar} &
    Orchestrator no portable \\
  \addlinespace
  Enterprise con SLA requerido &
    \textcolor{criticalred}{No adoptar} &
    Bus factor 1 inaceptable \\
  \addlinespace
  Stack Laravel o Angular &
    \textcolor{criticalred}{DIY manual} &
    Sin profile; wizard sin resultado util \\
  \addlinespace
  Pair programming puro &
    \textcolor{warnorange}{aider} &
    Mas directo, comunidad 50x mayor \\
  \addlinespace
  Autonomia a escala en nube &
    \textcolor{warnorange}{OpenHands} &
    Infraestructura que forge no tiene \\
  \addlinespace
  Flujo centrado en VS Code &
    \textcolor{warnorange}{cline / Roo Code} &
    Integracion nativa con el editor \\
  \bottomrule
\end{tabular}
\end{center}

%% ============================================================
\section{Conclusion}
%% ============================================================

\subsection{Para quien es forge ideal hoy}

Equipos de 2 a 8 personas que usan Claude Code de forma activa, trabajan con multiples
proyectos que comparten patrones de stack (Next.js, FastAPI, Expo, Rails), y necesitan
coherencia reproducible entre proyectos sin escribir instrucciones de agente desde cero
en cada uno. Especialmente para equipos con requisitos de compliance donde la propagacion
automatica de reglas tiene valor operativo real.

\subsection{Para quien no es forge adecuado hoy}

Desarrolladores individuales sin necesidad de coordinacion. Equipos que planean migrar
de Claude Code a otro runtime en el corto plazo. Equipos enterprise que necesitan un
vendor con soporte o SLA. Equipos con Laravel o Angular como stack principal.

\subsection{Veredicto}

forge v2.0 es el framework mas completo del ecosistema para gobernanza de agentes IA
en equipos que usan Claude Code. Su ventaja es real y documentada. Sus bugs tambien son
reales y documentados.

\textbf{\textcolor{criticalred}{El bug del campo \texttt{summary} en el JSON de CI
debe corregirse antes de cualquier adopcion en produccion:}} es una sola linea de codigo
con consecuencias operativas desproporcionadas. Corregido ese bug, forge es la eleccion
mas racional para el perfil de equipo descrito.

Para quienes no encajan en ese perfil, las alternativas son mejores en sus propios
dominios: aider para pair programming, OpenHands para autonomia a escala, cline para
integracion con VS Code. La eleccion no es ``forge vs. todo'': es ``forge para
gobernanza de equipos en Claude Code, las otras herramientas para todo lo demas''.

\subsection{Proximos pasos recomendados}

\begin{enumerate}
  \item \textbf{[Inmediato]} Corregir campo \texttt{summary} en \texttt{forge-audit.py}
    y actualizar ejemplos en README, guide.md y forge.py.
  \item \textbf{[Inmediato]} Implementar o eliminar flags \texttt{--forge} y
    \texttt{--only} de la documentacion.
  \item \textbf{[Corto plazo]} Actualizar documentacion a 13 profiles en forge.py,
    README y templates.
  \item \textbf{[Corto plazo]} Agregar tests del contrato JSON en
    \texttt{test\_forge\_audit.py}.
  \item \textbf{[Medio plazo]} Evaluar paquete pip o releases etiquetados como
    alternativa al submodule.
  \item \textbf{[Medio plazo]} Plan de gobernanza para reducir bus factor.
\end{enumerate}

\vspace{1cm}
\hrule
\vspace{0.5cm}
{\small\color{titlegray}
Documento generado a partir de analisis dual independiente sobre forge v2.0.
Fecha: 2026-05-03. No afiliado con Anthropic ni con ningun vendor mencionado en el benchmark.
}

\end{document}
